\section{The quaternionic Hecke estimate}\label{sec:exact-hecke}

The large-target arguments below use the same classical spectral input.
We state it once, allowing the fixed dyadic level needed for discriminant
$-8$.  The arithmetic sections identify their finite quotients, their
character components and their shell normalizations separately.

Let $B$ be a definite quaternion algebra over $\Q$, and write
$H=PB^\times$.  Let $U\subset H(\mathbb A_f)$ be an open compact
subgroup obtained from order units and congruence kernels.  Equip the
finite space
\[
 X=H(\Q)\backslash H(\mathbb A_f)/U
\]
with its invariant probability measure.  Before imposing a fine level,
the mass of a locally principal ideal class is proportional to the
reciprocal of the size of its projective unit group.  If the fine
stabilizers are trivial, this is uniform counting measure on fine
points, whose pushforward is the stated class measure.  It is essential
to retain every ideal class.

For an integer $m$ prime to the discriminant of $B$ and to the
nonmaximal or auxiliary level primes, let $T_m$ be the all-integral-ideal
Hecke operator, with maximal local level at every prime dividing $m$.
This includes imprimitive ideals.  Its constant eigenvalue is
$\sigma(m)$, and its local recurrence is
\[
 T_{\ell^{a+2}}=T_\ell T_{\ell^{a+1}}-\ell T_{\ell^a}.
\]
These conventions are those of
\cite[Propositions~41.3.1 and~41.3.6]{Voight}.

\begin{lemma}\label{lem:hecke-transfer}
The operators $T_m$ are commuting and self-adjoint for the invariant
measure.  On the orthogonal complement of all global reduced-norm
characters occurring in $L^2(X)$,
\[
                         \|T_m\|\le\tau(m)\sqrt m.
\]
More precisely, every simultaneous eigenvector in that complement has
unramified roots $\sqrt\ell\,u_\ell,\sqrt\ell\,v_\ell$ with
$|u_\ell|=|v_\ell|=1$ and $u_\ell v_\ell=1$.  Fixed operators at level
primes are not covered by this unramified assertion.
\end{lemma}

\begin{proof}
First, reversing an ideal flag gives the adjoint.  If $J\subset I$
has relative reduced norm $m$, elementary divisors at its split index
primes give $mI\subset J$.  The reversal
\[
                  (J\subset I)\longmapsto(mI\subset J)
\]
is an involution up to central scaling.  At auxiliary primes $m$ is a
unit and the local ideals are equal, so their markings transfer through
that equality; identifying $mI$ with $I$ changes the marking only by a
projectively trivial scalar.  Thus the reverse marked correspondence
has the same kernel.  On fine quotients it is symmetric, and on coarse
quotients its pushforward is symmetric for inverse-unit mass.  The
usual relatively prime ideal composition gives commutativity and the
displayed recurrence.

We specify the automorphic input and its hypotheses.  Decompose the
automorphic functions on this definite quotient, with trivial infinite
type, into irreducible constituents.  A finite-dimensional irreducible
smooth representation of $\mathrm{GL}_2(\Q_\ell)$ is a determinant
character.  Indeed its kernel is open, after intersecting the open
stabilizers of a basis; conjugation of elementary unipotents by diagonal
matrices and normality show that this kernel contains
$\mathrm{SL}_2(\Q_\ell)$.

If a constituent is a character at one split finite place $\ell$,
strong approximation for $B^1$, away from $\{\infty,\ell\}$, shows
that it is a global norm character.  To make the implication explicit,
for an automorphic function $f$, a finite adele $x$ and
$h\in B^1(\mathbb A_f)$, approximate $xhx^{-1}$ outside $\ell$ by
$\gamma\in B^1(\Q)$, sufficiently closely for the right stabilizers
of $f$.  Outside $\ell$, the points $xh$ and $\gamma x$ then have
the same value; at $\ell$ the discrepancy lies in the local norm-one
group, which acts trivially.  Left rational invariance gives
$f(xh)=f(x)$.  Thus the constituent factors through reduced norm.
The strong approximation hypothesis is satisfied because $B^1$ is
split and noncompact at $\ell$
\cite[Main Theorem~28.5.3]{Voight}.

After removing the global norm characters, every split local factor is
infinite dimensional.  Global Jacquet--Langlands therefore transfers
each remaining constituent to a cuspidal representation of
$\mathrm{GL}_2$ \cite[Theorem~14.4]{JL}.  Its trivial infinite type
transfers to weight two \cite[Section~5, Proposition~5.20]{JL}.
At each $\ell\mid m$ the quaternion algebra is split and the level
is maximal, so transfer preserves the unramified eigenvalues.  Fixed
levels elsewhere, including a nonmaximal dyadic level, do not change
this assertion.

The weight-two roots have absolute value $\sqrt\ell$: combine the
Frobenius--Hecke relation and purity in
\cite[Theorems~4.9 and~5.1]{DeligneModular} with
\cite[Theorem~1.6]{DeligneWeil}.  The central character is trivial,
so the normalized roots have product one.  The all-ideal eigenvalue
at $\ell^a$ is
\[
 \ell^{a/2}\sum_{j=0}^a u_\ell^jv_\ell^{a-j},
\]
whose absolute value is at most $(a+1)\ell^{a/2}$.
Multiplicativity and simultaneous self-adjointness give the claimed
operator norm.
\end{proof}
